% © 2026 Ing. David Jaroš — CC BY-NC-ND 4.0
%
% File: research_tracks/T3_ALPHA/electroweak_em_projection_theorem.tex
% Purpose: Reduce the final alpha proof condition by deriving the primitive
%          compact U(1)_EM projection and charge normalisation from the UBT
%          electroweak sector.
% Status: [L1 conditional] — proves the EM generator and unit-charge
%         normalisation conditional on the standard UBT/SM electroweak vacuum
%         representation. Does not derive the Weinberg angle.
% Date: 2026-06-13

\ifdefined\INCLUDEMODE\else
\documentclass[12pt,a4paper]{article}
\usepackage[a4paper,margin=2.5cm]{geometry}
\usepackage{amsmath,amssymb,amsthm,mathtools}
\usepackage{booktabs}
\usepackage{xcolor}
\usepackage[most]{tcolorbox}
\usepackage{hyperref}

\newtheorem{definition}{Definition}[section]
\newtheorem{lemma}[definition]{Lemma}
\newtheorem{theorem}[definition]{Theorem}
\newtheorem{corollary}[definition]{Corollary}
\newtheorem{remark}[definition]{Remark}

\newtcolorbox{statusbox}[1][]{
  colback=yellow!8!white,
  colframe=orange!70!black,
  arc=3pt,
  boxrule=1pt,
  title=Status,
  #1
}
\newtcolorbox{resultbox}[1][]{
  colback=green!8!white,
  colframe=green!50!black,
  arc=3pt,
  boxrule=1pt,
  title=Result,
  #1
}

\newcommand{\Lone}{\textbf{[L1 cond.]}}
\newcommand{\Std}{\textbf{[STD]}}

\title{Electroweak Projection to Primitive Compact $U(1)_{\rm EM}$\\
       \large and Charge Normalisation for the UBT Alpha Track}
\author{Ing.\ David Jaro\v{s}}
\date{2026-06-13}

\begin{document}
\maketitle
\fi

\begin{abstract}
The alpha derivation requires the UBT winding modulus to be the compact
unit-charge electromagnetic coupling modulus.  This note isolates and proves
the corresponding electroweak projection statement at conditional level.  The
UBT electroweak sector contains \(SU(2)_L\) from the left quaternionic action
and \(U(1)_Y\) from the right scalar phase.  If the low-energy vacuum transforms
as the standard electroweak doublet with hypercharge \(Y=1\), its stabiliser is
the primitive compact generator
\[
  Q_{\rm EM}=T_3+\frac{Y}{2}.
\]
The primitive Wilson-line normalisation fixes the unit charge and forbids an
extra continuous rescaling of \(Q_{\rm EM}\).  This proves the projection needed
by the alpha EM modular identification.  It does not derive the Weinberg angle;
the angle controls the decomposition of the photon field into \(W^3\) and \(B\),
not the primitive electromagnetic charge generator.
\end{abstract}

\begin{statusbox}
\textbf{Classification.} Conditional theorem.  Inputs already present in the
UBT gauge sector are used: \(SU(2)_L\) from left action and \(U(1)_Y\) from
right phase.  The remaining physical condition is the vacuum representation:
\(\Theta_0\) must be the electroweak doublet with \(Y=1\).  The theorem proves
that, under that condition, the residual unbroken gauge group is the primitive
compact \(U(1)_{\rm EM}\) with no free charge-normalisation factor.
\end{statusbox}

\section{Electroweak action on the UBT doublet}

Let \(\Theta\) contain the low-energy electroweak doublet component
\[
  \Phi=\begin{pmatrix}\phi^+\\ \phi^0\end{pmatrix},
\]
with \(SU(2)_L\) generators
\[
  T_a=\frac{\sigma_a}{2},
  \qquad
  T_3=\frac12\begin{pmatrix}1&0\\0&-1\end{pmatrix},
\]
and hypercharge generator \(Y\) acting by the right scalar phase.  In the
standard electroweak convention the Higgs-like vacuum has
\[
  \Theta_0\supset \Phi_0=\begin{pmatrix}0\\v/\sqrt2\end{pmatrix},
  \qquad
  Y\Phi_0=1\cdot\Phi_0 .
\]
Thus
\[
  T_3\Phi_0=-\frac12\Phi_0,
  \qquad
  \frac{Y}{2}\Phi_0=+\frac12\Phi_0.
\]

\section{Stabiliser theorem}

\begin{theorem}[Residual electromagnetic generator \Lone]
If the UBT low-energy vacuum contains the electroweak doublet component
\(\Phi_0=(0,v/\sqrt2)^T\) with hypercharge \(Y=1\), then the unique unbroken
one-dimensional generator in \(\mathfrak{su}(2)_L\oplus\mathfrak{u}(1)_Y\) is
\[
  Q_{\rm EM}=T_3+\frac{Y}{2}.
\]
\end{theorem}

\begin{proof}
Let a general neutral generator be
\[
  X=aT_3+b\frac{Y}{2}.
\]
The vacuum is invariant iff \(X\Phi_0=0\).  Since
\[
  T_3\Phi_0=-\frac12\Phi_0,
  \qquad
  \frac{Y}{2}\Phi_0=+\frac12\Phi_0,
\]
we get
\[
  X\Phi_0=\frac{-a+b}{2}\Phi_0.
\]
Therefore \(X\Phi_0=0\) iff \(a=b\).  Up to an overall normalisation, the
unique unbroken generator is
\[
  Q_{\rm EM}=T_3+\frac{Y}{2}.
\]
The two orthogonal generators \(T_1,T_2\) and the neutral orthogonal combination
\(T_3-Y/2\) are broken by the vacuum.
\end{proof}

\section{Primitive compact normalisation}

\begin{definition}[Primitive electromagnetic generator]
A compact \(U(1)\) generator \(Q\) is primitive if the Wilson loop
\[
  \exp(2\pi i Q)
\]
acts trivially on all allowed physical fields, and no smaller non-zero
rescaling of \(Q\) has this property while preserving the observed unit charge.
\end{definition}

\begin{theorem}[No continuous charge rescaling \Lone]
The generator
\[
  Q_{\rm EM}=T_3+\frac{Y}{2}
\]
with the electron assigned charge \(-1\) is primitive.  Therefore the
low-energy electromagnetic line bundle is a compact unit-charge
\(U(1)_{\rm EM}\) bundle.  There is no continuous free factor
\(Q_{\rm EM}\mapsto\lambda Q_{\rm EM}\).
\end{theorem}

\begin{proof}
The charged lepton doublet has charges
\[
  Q(\nu_L)=0,
  \qquad
  Q(e_L)=-1,
\]
and the right electron has \(Q(e_R)=-1\).  Thus the smallest observed lepton
charge magnitude is one in this normalisation.  A rescaling
\(Q\mapsto\lambda Q\) changes the Wilson-line phases
\(\exp(2\pi i Q)\) and changes the electron charge unless \(\lambda=1\) within
the chosen unit-charge convention.  Fractional quark charges are compatible
with the usual colour/electroweak quotient structure, but they do not change
the primitive electromagnetic unit fixed by the electron Wilson line.
Therefore the residual \(U(1)\) is compact and unit-charge normalised.
\end{proof}

\section{Relation to the Weinberg angle}

The physical photon connection may be written
\[
  A_\mu=\sin\theta_W\,W^3_\mu+\cos\theta_W\,B_\mu,
\]
with
\[
  e=g\sin\theta_W=g'\cos\theta_W.
\]
This angle determines how the unbroken \(U(1)_{\rm EM}\) connection is embedded
inside the original \(SU(2)_L\times U(1)_Y\) gauge potentials.  It does not
change the primitive generator
\[
  Q_{\rm EM}=T_3+\frac{Y}{2},
\]
and it does not introduce a free rescaling of the compact electromagnetic
Wilson-line unit.  Therefore the alpha track may compute the final low-energy
\(U(1)_{\rm EM}\) coupling modulus directly; it need not derive
\(\sin^2\theta_W\) as a separate route to \(\alpha\).

\begin{remark}
This is compatible with the existing UBT gauge-status statement that the pure
algebra of \(\mathbb C\otimes\mathbb H\) does not fix \(g'/g\).  The present
result does not claim to derive the Weinberg angle.  It only proves the
unbroken compact electromagnetic generator and its primitive charge
normalisation, which is the part needed by the alpha EM modular-coupling
identification.
\end{remark}

\section{Connection to the alpha theorem}

Combining this projection theorem with the EM modular identification gives
\[
  \tau^{\rm UBT}_{\rm EM}=\chi+i n
  \quad\text{and}\quad
  \tau_{\rm EM}=\frac{\theta_{\rm EM}}{2\pi}+i\frac{4\pi}{e^2}.
\]
At the CP-even vacuum \(\chi=\theta_{\rm EM}/(2\pi)=0\),
\[
  n=\frac{4\pi}{e^2}=\alpha^{-1}.
\]
Thus the remaining alpha proof condition is reduced to the vacuum theorem:
\[
  \boxed{
    \Theta_0\text{ contains the standard }(SU(2)_L,Y=1)\text{ doublet vacuum.}
  }
\]

\begin{resultbox}
\textbf{Result.} The EM projection needed by the alpha track is no longer an
arbitrary identification.  Under the standard electroweak-vacuum condition, the
residual group is the primitive compact unit-charge \(U(1)_{\rm EM}\).  This
removes the possible free normalisation factor in \(n=\alpha^{-1}\).
\end{resultbox}

\ifdefined\INCLUDEMODE\else
\end{document}
\fi
